%! Minimizing Loss by Gradient Descent — Unit 8, Lesson 8.3 — Group Activity
\documentclass[10pt]{article}
\usepackage{linalg-article}
\usepackage{linalg-boxes}
\newcommand{\TallMath}[1]{$\displaystyle #1\rule[-1.4em]{0pt}{3.2em}$}
\newcommand{\ww}{\mathbf{w}}
\newcommand{\zero}{\mathbf{0}}

\begin{document}

\pageheader{Unit 8, Lesson 8.3}{Group Activity}
\namepartnerperiod

\vspace{0.10in}

\begin{headlinebox}{blush}
\small\textbf{Rolling Downhill.} \emph{Gradient descent} chooses a network's weights by rolling downhill on the
loss: at each step, feel the slope and step the \emph{opposite} way, $w\leftarrow w-s\,L'(w)$. The learning
rate $s$ is the step size; the bottom of the bowl (slope $=0$) is the best weight. Work with your partner, show
every step, and \emph{explain} what you find.
\end{headlinebox}

\vspace{0.10in}

\begin{tcolorbox}[colback=white, colframe=black!40, title=\textbf{Tier R --- Remediate},
    fonttitle=\bfseries, arc=1.5mm, left=3mm, right=3mm, top=2mm, bottom=2mm, breakable]
Use the bowl $L(w)=(w-4)^2$ with slope $L'(w)=2w-8$ and learning rate $s=0.25$.
\begin{enumerate}[leftmargin=1.7em, itemsep=8pt]
  \item \textbf{Two steps from $w=0$.} Fill the table (slope, then $w-s\,L'(w)$):
        \[ \begin{array}{c|c|c} w & L'(w)=2w-8 & \text{new } w=w-0.25\,L'(w)\\ \hline
            0 & \blank{1.0cm} & \blank{1.2cm}\\[2pt]
            2 & \blank{1.0cm} & \blank{1.2cm}\\ \end{array} \]
  \item \textbf{Watch the loss fall.} Compute the loss $L(w)=(w-4)^2$ at your three weights:
        \[ L(0)=\blank{0.9cm},\qquad L(2)=\blank{0.9cm},\qquad L(3)=\blank{0.9cm}. \]
        In one sentence, what is happening to the loss as $w$ moves? \writeline
\end{enumerate}
\end{tcolorbox}

\begin{tcolorbox}[colback=white, colframe=black!40, title=\textbf{Tier A --- Approaching Proficiency},
    fonttitle=\bfseries, arc=1.5mm, left=3mm, right=3mm, top=2mm, bottom=2mm, breakable]
Same bowl $L(w)=(w-4)^2$, slope $L'(w)=2w-8$.
\begin{enumerate}[leftmargin=1.7em, itemsep=9pt]
  \item \textbf{Descend from the other side.} Start at $w=6$ with $s=0.25$ and take two steps.
        \[ w:\ 6\ \to\ \blank{1.0cm}\ \to\ \blank{1.0cm}. \]
        Which direction did $w$ move this time, and toward what? \writeline
  \item \textbf{Know when to stop.} Compute the slope at the minimum, $L'(4)$. What does the update
        $w\leftarrow w-s\,L'(w)$ do when the slope is this value, and why is that the signal that you have
        arrived? \writeline
  \item \textbf{The learning rate.} Take \emph{one} step from $w=0$ with each rate and compare:
        \[ s=0.5:\ w=0-0.5(-8)=\blank{1.0cm},\qquad s=1:\ w=0-1(-8)=\blank{1.0cm}. \]
        One lands exactly at the minimum; the other overshoots. In a sentence, what does too big a learning
        rate do? \writeline
\end{enumerate}
\end{tcolorbox}

\begin{tcolorbox}[colback=white, colframe=black!40, title=\textbf{Tier E --- Extension},
    fonttitle=\bfseries, arc=1.5mm, left=3mm, right=3mm, top=2mm, bottom=2mm, breakable]
\textbf{Two weights at once --- the gradient.} For $L(w_1,w_2)=(w_1-1)^2+(w_2-3)^2$ the gradient is
$\nabla L=[\,2(w_1-1),\,2(w_2-3)\,]$; the rule is $\ww\leftarrow\ww-s\,\nabla L$ with $s=0.25$.
\begin{enumerate}[leftmargin=1.7em, itemsep=9pt]
  \item \textbf{Two steps from $(0,0)$.}
        \[ \nabla L(0,0)=[\ \blank{0.6cm},\ \blank{0.6cm}\ ],\quad (0,0)-0.25\,\nabla L=(\ \blank{0.6cm},\ \blank{0.6cm}\ ); \]
        \[ \text{then from }(0.5,1.5):\ \nabla L=[\ \blank{0.6cm},\ \blank{0.6cm}\ ],\quad \text{new }\ww=(\ \blank{0.6cm},\ \blank{0.6cm}\ ). \]
        Toward which point is $\ww$ heading? \writeline
  \item \textbf{Justify.} In one or two sentences each: (a) why do we \emph{subtract} $s\,\nabla L$ (step
        \emph{against} the gradient) instead of adding it? (b) Why does slope $=0$ (a flat gradient) mean we
        have found the best weights? \writeline
\end{enumerate}
\end{tcolorbox}

\end{document}
